\subsection{The direct topological sum of a family of locally convex spaces}

DEFINITION 2. --- Let E be the vector space which is the direct sum (A, II, §1.6) of the family of locally convex spaces \((\mathrm{E}_{\iota})_{\iota \in \mathbf{I}}\) . For each \(\iota \in \mathbf{I}\) , let \(f_{\iota}\) be the canonical injection of \(E_{t}\) in E. By the topological direct sum of the family ( \(E_{t}\) ) we mean the space E with the finest locally convex topology which makes each \(f_{t}\) continuous (this topology is called the direct sum of the topologies of the \(E_{t}\) ).

In the remainder of this No.~we keep the same notations as in def. 2 (unless the contrary is expressly stated) and we identify, canonically, each \(E_{t}\) with a subspace of E, by means of \(f_{t}\) .

By the general description of neighbourhoods of a locally convex final topology given in II, p.~28, we can here obtain a fundamental system of neighbourhoods of 0 in E for the direct sum topology, in the following manner; for every family \((\mathrm{V}_{\iota})_{\iota \in I}\) where \(\mathrm{V}_{\iota}\) is a symmetric convex neighbourhood of 0 in \(\mathrm{E}_{\iota}\) , consider the convex envelope \(\Gamma((\mathrm{V}_{\iota}))\) , of the union of the \(\mathrm{V}_{\iota}\) ; the \(\Gamma((\mathrm{V}_{\iota}))\) for all the families \((\mathrm{V}_{\iota})\) (or only taking \(\mathrm{V}_{\iota}\) for each \(\iota\) to be in a fundamental system of neighbourhoods of 0 in \(\mathrm{E}_{\iota}\) ) form a fundamental system of neighbourhoods of 0 in E.

Example. --- Let \((a_{\mathfrak{t}})_{\mathfrak{t}\in\mathbf{I}}\) be a basis of the vector space E and consider the canonical topology (I, p.~2, Example 5) on each line \(Ra_{\mathfrak{t}}\) ; the direct sum of these topologies is the finest locally convex topology on E (II, p.~26); in fact, if V is an absorbent, symmetric, convex set in E, then \(V_{\mathfrak{t}} = V \cap Ra_{\mathfrak{t}}\) is a neighbourhood of 0 in \(Ra_{\mathfrak{t}}\) and V clearly contains the convex envelope \(\Gamma((V_{\mathfrak{t}}))\) .

PROPOSITION 6. --- A locally convex topology T on E is the direct sum of the topologies of the \(E_{t}\) , if and only if, the following property holds: a linear mapping of E in a locally convex space G (resp. a semi-norm p on E) is continuous, if and only if, for every \(t \in I\) , the mapping \(g \circ f_{t}\) (resp. \(p \circ f_{t}\) ) is continuous in \(E_{t}\) .

This is a particular case of prop. 5, II, p.~27.

Recalling the definition of the direct sum of a family of vector spaces (A, II, p.~12, prop. 6), we can say that the topology \(\mathcal{T}\) is the only one for which the canonical mapping \(g \mapsto (g \circ f_{\mathrm{t}})\) is a bijection

\[
\mathcal {L} (\mathrm{E}; \mathrm{G}) \rightarrow \prod_ {\iota \in \mathrm{I}} \mathcal {L} (\mathrm{E} _ {\iota}; \mathrm{G})\tag{1}
\]

for every locally convex space \(\mathbf{G}\) .

COROLLARY. --- With the notation of prop. 5, II, p.~27, suppose that E is the sum of the \(g_{\alpha}(\mathbf{F}_{\alpha})\) . Let F be the topological direct sum of the family \((\mathbf{F}_{\alpha})_{\alpha \in \mathbf{A}}\) , and let \(j_{\alpha}: \mathbf{F}_{\alpha} \to \mathbf{F}\) be the canonical injection; suppose that \(g: \mathbf{F} \to \mathbf{E}\) is the linear mapping such that \(g \circ j_{\alpha} = g_{\alpha}\) for all \(\alpha \in \mathbf{A}\) . If N is the kernel of \(g\) , then the canonical bijection \(\mathrm{F/N} \to \mathrm{E}\) associated with \(g\) is a topological isomorphism of \(\mathrm{F/N}\) on E with the topology \(\mathcal{T}\) .

This is a particular case of II, p.~28, cor. 2 remembering II, p.~29, Example I.

PROPOSITION 7. --- The canonical injection \(j: \mathrm{E} \to \prod_{\iota \in \mathrm{I}} \mathrm{E}_{\iota}\) is continuous when \(\mathrm{E}\) carries the direct sum topology of the \(\mathrm{E}_{\iota}\) and \(\prod_{\iota \in \mathrm{I}} \mathrm{E}_{\iota}\) carries the product topology. When \(\mathrm{I}\) is finite, this mapping is an isomorphism of topological vector spaces.

The first assertion follows from the fact that the canonical injections \(E_{\kappa} \to \prod_{i \in I} E_{i}\) are continuous for each \(\kappa \in I\) . If I is finite then j is the identity mapping, and it is sufficient to show that the product topology \(T'\) is finer than the direct sum topology T. Now, let V be a convex neighbourhood of 0 for T; each set \(V \cap E_{i}\) is a convex neighbourhood of 0 in \(E_{i}\) ; if n is the number of elements of I, then the set V contains the set \(\frac{1}{n} \sum_{n} (V \cap E_{i})\) , which is a neighbourhood of 0 for \(T'\) , and the proposition is proved.

When I is infinite, if, for each finite subset J of I, we write \(E_{J}\) for the space \(\prod_{i\in J}E_{i}\) , with the product topology, then E is the inductive limit of the \(E_{J}\) (identified as subspaces of E).

PROPOSITION 8. --- Let \(\mathbf{N}_{\iota}\) be a subspace of \(\mathrm{E}_{\iota}\) , for every \(\iota \in \mathbf{I}\) ,

\begin{enumerate}
\def\labelenumi{(\roman{enumi})}
\item
  The topology induced on \(\mathbf{N} = \sum_{\mathfrak{i}}\mathbf{N}_{\mathfrak{i}}\) by the direct sum topology \(\mathcal{T}\) on \(\mathbf{E}\) is identical with the direct sum of the topologies of the \(\mathbf{N}_{\mathfrak{i}}\) .
\item
  The canonical mapping \(h\) of the topological direct sum space of the \(\mathrm{E}_{\mathrm{i}} / \mathrm{N}_{\mathrm{i}}\) on \(\mathrm{E} / \mathrm{N}\) (A, II, § 1.6, formula (26)) is an isomorphism between topological vector spaces.
\item
  With the notations introduced above, we consider \(x = \sum_{i} \lambda_i x_i\) belonging to \(N \cap \Gamma((V_{t}))\) ( \((\lambda_{t})\) is a family of numbers \(\geqslant 0\) of which at most finitely many are non-zero, such that \(\sum_{i} \lambda_{i} = 1\) , and \(x_{i} \in V_{i}\) , for all \(i \in I\) ).
\end{enumerate}

Since the sum of the \(N_{i}\) is direct, we have \(\lambda_{i}x_{i} \in N_{i}\) for all \(i \in I\) ; therefore, for all i such that \(\lambda_{i} > 0\) we also have \(x_{i} \in N_{i} \cap V_{i}\) , and x belongs to the convex envelope \(\Gamma((N_{i} \cap V_{i}))\) , thus (i) is proved. (ii) Denote canonical mappings as follows: \(f_{i}: E_{i} \to E\) , \(h_{i}: E_{i}/N_{i} \to E/N\) , \(p_{i}: E_{i} \to E_{i}/N_{i}\) and \(p: E \to E/N\) . For every \(i \in I\) , \(h_{i} \circ p_{i} = p \circ f_{i}\) and the proposition follows from II, p.~28, cor. 2 and p.~29 Example I.

COROLLARY 1. --- If \(\mathbf{N}_{\iota}\) is closed in \(\mathbf{E}_{\iota}\) for every \(\iota \in \mathbf{I}\) , then \(\mathbf{N} = \sum \mathbf{N}_{\iota}\) is closed in \(\mathbf{E}\) .

For, the canonical mapping \(p_{\iota} : \mathrm{E} \to \mathrm{E}_{\iota}\) is continuous (II, § 4.5, prop. 6) for every \(\iota \in \mathbf{I}\) , hence \(p_{\iota}^{-1}(\mathbf{N}_{\iota})\) is closed in \(\mathbf{E}\) , and thus the same is true of the intersection \(\mathbf{N} = \bigcap_{\iota \in \mathbf{I}} p_{\iota}^{-1}(\mathbf{N}_{\iota})\) .

COROLLARY 2. --- If each \(E_{t}\) is Hausdorff, so also is E and each \(E_{t}\) is closed in E.

To prove the first statement apply cor. 1 taking \(N_{t} = \{0\}\) for all \(t \in I\) ; for the second apply cor. 1 with \(N_{t} = E_{t}\) and \(N_{\kappa} = \{0\}\) for every \(\kappa \neq t\) .

We shall show in III, p.~21, cor. 2 that if the \(E_{t}\) are Hausdorff and complete then so is their topological direct sum E.

